\documentclass{article}
\usepackage[T1]{fontenc}
\usepackage{lmodern}
\usepackage{graphicx}
\usepackage{amsmath,amssymb}
\usepackage[a4paper,margin=2cm]{geometry}
\usepackage[absolute,overlay]{textpos}
\setlength{\TPHorizModule}{1mm}
\setlength{\TPVertModule}{1mm}
\usepackage[indonesian]{babel}
\usepackage{hyperref}
\setlength{\emergencystretch}{2em}
% unit: Halaman 19

\begin{document}

% === Halaman 19 (llm) ===

\section*{Pertukaran kunci Diffie-Hellman untuk tiga pihak (Alice, Bob, dan Carol)}

\begin{enumerate}
  \item Alice, Bob, dan Carol menyepakati $p$ dan $g$.
  \item Alice, Bob, dan Carol membangkitkan kunci privat masing-masing, $a$, $b$, dan $c$.
  \item Alice menghitung $g^a \pmod p$ dan mengirimkannya kepada Bob.
  \item Bob menghitung $(g^a)^b \pmod p = g^{ab} \pmod p$ dan mengirimkannya kepada Carol.
  \item Carol menghitung $K = (g^{ab})^c \pmod p = g^{abc} \pmod p$.
  \item Bob menghitung $g^b \pmod p$ dan mengirimkannya kepada Carol.
  \item Carol menghitung $(g^b)^c \pmod p = g^{bc} \pmod p$ dan mengirimkannya kepada Alice.
  \item Alice menghitung $K = (g^{bc})^a \pmod p = g^{bca} \pmod p = g^{abc} \pmod p$.
  \item Carol menghitung $g^c \pmod p$ dan mengirimkannya kepada Alice.
  \item Alice menghitung $(g^c)^a \pmod p = g^{ca} \pmod p$ dan mengirimkannya kepada Bob.
  \item Bob menghitung $K = (g^{ca})^b \pmod p = g^{cab} \pmod p = g^{abc} \pmod p$.
  \item Sekarang Alice, Bob, dan Carol sudah memiliki kunci rahasia yang sama, yaitu $K$
\end{enumerate}

\end{document}
